\chapter{Pancreaticoduodenectomy (Whipple Procedure)}
\section*{Overview}
The Whipple procedure removes the pancreatic head, duodenum, gallbladder, and part of the bile duct, followed by reconstruction of the digestive tract.
\section*{Why It Is Done}
It is most often performed for resectable pancreatic-head cancer and selected tumors or diseases of the bile duct, duodenum, or ampulla.
\section*{How It Is Performed}
The involved organs are removed and the remaining pancreas, bile duct, and stomach or intestine are connected to the small bowel to restore drainage and digestion.
\section*{Risks and Recovery}
Pancreatic leak, bleeding, infection, delayed gastric emptying, diabetes, digestive problems, and mortality are possible. Recovery is prolonged and requires nutritional and oncology follow-up.
\section*{References}
\begin{enumerate}\item National Cancer Institute. Pancreatic Cancer Surgery. 2025.\item Current Medical Diagnosis and Treatment. McGraw-Hill; 2025.\end{enumerate}
\clearpage
